% ENEM 2015-2025 — 7 questões MCQ — Aritmética III
% Fonte: lista "Aritmética III" do site projetoagathaedu.com.br
% Omitidas: Q1 (duplicata de Q4 com gabarito inconsistente), Q3,Q8,Q9,Q11,Q12 (imagens)
%           Q10 (duplicata idêntica de Q6)

---
tipo: Múltipla Escolha
dificuldade: Fácil
nivel: médio
disciplina: Matemática
assunto: Aritmética
gabarito: A
tags: [ENEM, aritmética, custo, estacionamento, inequação]
fonte: concurso
concurso: ENEM
banca: INEP
ano: 2016
numero: 2
---
\question
O gerente de um estacionamento, próximo a um grande aeroporto, sabe que um passageiro que utiliza seu carro nos traslados casa-aeroporto-casa gasta cerca de R\$ 10,00 em combustível nesse trajeto. Ele sabe, também, que um passageiro que não utiliza seu carro nos traslados casa-aeroporto-casa gasta cerca de R\$ 80,00 com transporte. Suponha que os passageiros que utilizam seus próprios veículos deixem seus carros nesse estacionamento por um período de dois dias.

Para tornar atrativo a esses passageiros o uso do estacionamento, o valor, em real, cobrado por dia de estacionamento deve ser, no máximo, de
\begin{choices}
\CorrectChoice 35,00.
\choice 40,00.
\choice 45,00.
\choice 70,00.
\choice 90,00.
\end{choices}

---
tipo: Múltipla Escolha
dificuldade: Fácil
nivel: médio
disciplina: Matemática
assunto: Progressão Geométrica
gabarito: C
tags: [ENEM, PG, progressão geométrica, visitantes, triplicar]
fonte: concurso
concurso: ENEM
banca: INEP
ano: 2016
numero: 4
---
\question
Para comemorar o aniversário de uma cidade, a prefeitura organiza quatro dias consecutivos de atrações culturais. A experiência de anos anteriores mostra que, de um dia para o outro, o número de visitantes no evento é triplicado. É esperada a presença de 345 visitantes para o primeiro dia do evento.

Uma representação possível do número esperado de participantes para o último dia é
\begin{choices}
\choice $3 \times 345$
\choice $(3 + 3 + 3) \times 345$
\CorrectChoice $3^3 \times 345$
\choice $3 \times 4 \times 345$
\choice $3^4 \times 345$
\end{choices}

---
tipo: Múltipla Escolha
dificuldade: Média
nivel: médio
disciplina: Matemática
assunto: Grandezas e Medidas
gabarito: B
tags: [ENEM, grandezas, volume, plasma, banco de sangue, congelador]
fonte: concurso
concurso: ENEM
banca: INEP
ano: 2016
numero: 5
---
\question
Um banco de sangue recebe 450 mL de sangue de cada doador. Após separar o plasma sanguíneo das hemácias, o primeiro é armazenado em bolsas de 250 mL de capacidade. O banco de sangue aluga refrigeradores de uma empresa para estocagem das bolsas de plasma, segundo a sua necessidade. Cada refrigerador tem uma capacidade de estocagem de 50 bolsas. Ao longo de uma semana, 100 pessoas doaram sangue àquele banco.

Admita que, de cada 60 mL de sangue, extraem-se 40 mL de plasma.

O número mínimo de congeladores que o banco precisou alugar, para estocar todas as bolsas de plasma dessa semana, foi
\begin{choices}
\choice 2.
\CorrectChoice 3.
\choice 4.
\choice 6.
\choice 8.
\end{choices}

---
tipo: Múltipla Escolha
dificuldade: Fácil
nivel: médio
disciplina: Matemática
assunto: Grandezas e Medidas
gabarito: C
tags: [ENEM, grandezas, dosagem, antibiótico, coelhos, volume]
fonte: concurso
concurso: ENEM
banca: INEP
ano: 2015
numero: 6
---
\question
Um granjeiro detectou uma infecção bacteriológica em sua criação de 100 coelhos. A massa de cada coelho era de, aproximadamente, 4 kg. Um veterinário prescreveu a aplicação de um antibiótico, vendido em frascos contendo 16 mL, 25 mL, 100 mL, 400 mL ou 1.600 mL. A bula do antibiótico recomenda que, em aves e coelhos, seja administrada uma dose única de 0,25 mL para cada quilograma de massa do animal.

Para que todos os coelhos recebessem a dosagem do antibiótico recomendada pela bula, de tal maneira que não sobrasse produto na embalagem, o criador deveria comprar um único frasco com a quantidade, em mililitros, igual a
\begin{choices}
\choice 16.
\choice 25.
\CorrectChoice 100.
\choice 400.
\choice 1.600.
\end{choices}

---
tipo: Múltipla Escolha
dificuldade: Fácil
nivel: médio
disciplina: Matemática
assunto: Números Decimais
gabarito: C
tags: [ENEM, números decimais, comparação, lentes, espessura]
fonte: concurso
concurso: ENEM
banca: INEP
ano: 2025
numero: 7
---
\question
Deseja-se comprar lentes para óculos. As lentes devem ter espessuras mais próximas possíveis da medida 3 mm. No estoque de uma loja, há lentes de espessuras: 3,10 mm; 3,021 mm; 2,96 mm; 2,099 mm e 3,07 mm.

Se as lentes forem adquiridas nessa loja, a espessura escolhida será, em milímetros, de
\begin{choices}
\choice 2,099.
\choice 2,96.
\CorrectChoice 3,021.
\choice 3,07.
\choice 3,10.
\end{choices}

---
tipo: Múltipla Escolha
dificuldade: Média
nivel: médio
disciplina: Matemática
assunto: Divisibilidade
gabarito: E
tags: [ENEM, MDC, divisibilidade, tábuas, carpinteiro, corte]
fonte: concurso
concurso: ENEM
banca: INEP
ano: 2025
numero: 13
---
\question
Um arquiteto está reformando uma casa. De modo a contribuir com o meio ambiente, decide reaproveitar tábuas de madeira retiradas da casa. Ele dispõe de 40 tábuas de 540 cm, 30 de 810 cm e 10 de 1.080 cm, todas de mesma largura e espessura. Ele pediu a um carpinteiro que cortasse as tábuas em pedaços de mesmo comprimento, sem deixar sobras, e de modo que as novas peças ficassem com o maior tamanho possível, mas de comprimento menor que 2 m.

Atendendo o pedido do arquiteto, o carpinteiro deverá produzir
\begin{choices}
\choice 105 peças.
\choice 120 peças.
\choice 210 peças.
\choice 243 peças.
\CorrectChoice 420 peças.
\end{choices}

---
tipo: Múltipla Escolha
dificuldade: Média
nivel: médio
disciplina: Matemática
assunto: Divisibilidade
gabarito: C
tags: [ENEM, MDC, divisibilidade, ingressos, escola, distribuição]
fonte: concurso
concurso: ENEM
banca: INEP
ano: 2025
numero: 14
---
\question
O gerente de um cinema fornece anualmente ingressos gratuitos para escolas. Este ano serão distribuídos 400 ingressos para uma sessão vespertina e 320 ingressos para uma sessão noturna de um mesmo filme. Há alguns critérios para a distribuição dos ingressos: 1) cada escola deverá receber ingressos para uma única sessão; 2) todas as escolas contempladas deverão receber o mesmo número de ingressos; 3) não haverá sobra de ingressos (ou seja, todos os ingressos serão distribuídos).

O número mínimo de escolas que podem ser escolhidas para obter ingressos, segundo os critérios estabelecidos, é
\begin{choices}
\choice 2.
\choice 4.
\CorrectChoice 9.
\choice 40.
\choice 80.
\end{choices}
